\section{Conclusion}
Representing a discrete-time signal requires both its sample vector and
its time origin. Time shifts change the origin, folding reverses the
sample order and updates the origin, and addition or pointwise
multiplication requires alignment by time. Convolution accumulates
products at sums of input time indices and has a different origin rule.

The supplementary derivations establish that folding and shifting do
not generally commute and that the decomposition into even and odd
components exists and is unique. For the three systems in Problem 2.6,
the first-difference system is time invariant, while the squared-time
and time-multiplication systems are time varying. A failed shift test
provides a counterexample; proving time invariance requires an identity
valid for arbitrary inputs and shifts.
